\documentclass{article}

\usepackage[utf8]{inputenc}
\usepackage[T1]{fontenc}
\usepackage[english]{babel}
\usepackage{amsmath}
\usepackage{amssymb}
\usepackage{booktabs}
\usepackage{hyperref}
\usepackage[a4paper,margin=2.5cm]{geometry}

\title{The Snake Game\\\large Modern Computational Physics I}
\author{Ján Kovačovský}
\date{}

\begin{document}

\maketitle

\section{Problem formulation}

The project studies a deterministic variant of the Snake game on a fixed $20 \times 20$ grid with wrap-around boundaries. Leaving the grid on one side immediately places the snake on the opposite side, so the playing field is topologically a torus. The initial snake head position and the set of apples are fixed for all experiments, which makes the benchmark directly comparable across different methods.

The task is to eat all apples while minimizing the total trajectory length. Since the apples are static, the problem is closely related to a travelling-salesman-style ordering problem on a wrapped grid. Unlike a pure travelling salesman formulation, however, the snake body grows after each collected apple, so a route is only valid if it also avoids self-collisions during execution.

This makes it a more challenging combinatorial optimization problem than a standard TSP, and it also means that the optimal solution is not necessarily the shortest path through the apple locations. The benchmark instance is small enough to allow for an effectively optimal solution by A* search, which serves as a reference for evaluating the quality of approximate and learned methods.

\section{Approaches}

We compare three broad families of methods.

\subsection{A* search}

A* serves as the reference method. It searches directly in the discrete state space of the game, where a state is determined by the current snake head position together with the set of apples that still remain to be collected. Because the benchmark instance is fixed and moderately sized, A* is able to provide an effectively optimal solution for the chosen board. This makes it a natural baseline against which all approximate and learned methods can be measured.

\subsection{Route optimization methods}

The second family consists of route-based optimizers. These methods treat the problem primarily as an ordering of apple visits and optimize a cost function derived from the resulting path length. Three variants were considered:

\begin{itemize}
    \item \textbf{Genetic algorithm}, which evolves a population of candidate visit orders by mutation and selection,
    \item \textbf{simulated annealing}, which performs stochastic local search with a temperature schedule that gradually reduces exploration, and
    \item \textbf{Metropolis sampling}, which uses related stochastic proposals but without the full annealing schedule.
\end{itemize}

These approaches are attractive because they are simple to formulate and can explore a large combinatorial search space without exhaustive enumeration. Their disadvantage is that they optimize only an approximate route representation, so the final quality depends on how well the optimized order translates into a collision-free snake trajectory.

\subsection{Learned function approximators}

The third family contains learned models that attempt to represent a policy or value landscape for the game and then improve their parameters by numerical optimization. Two model classes were used:

\begin{itemize}
    \item a small deep Q-network (DQN) implemented as a compact multilayer perceptron, and
    \item a Kohonen-style model used as an alternative lightweight learned representation.
\end{itemize}

For these learned solvers, the main comparison is not only between model families, but also between optimization algorithms used during training. In addition to the standard \texttt{Adam} optimizer from PyTorch, we tested optimizers from the \texttt{optimizers} submodule, in particular the \textbf{Extended Kalman Filter} (EKF) and \textbf{Levenberg--Marquardt} (LM) methods. The motivation is that these second-order or quasi-second-order methods can converge faster or reach better minima on small neural models than first-order gradient descent alone.

\section{Common experimental setup}

All approaches were run on the same wrapped $20 \times 20$ board with the same initial head position and the same 13 apples. This eliminates randomness in the problem instance itself and ensures that differences in trajectory quality and runtime come from the compared methods rather than from different game layouts.

The backend exports each run into a common replay JSON format. This replay file is the canonical output of the experiment because it captures the actual executed trajectory, the eaten apples, and the frame-by-frame movement used by the visualization. Consequently, the main quality metric is the replay length, with the A* replay serving as the optimal reference on this board.

To make runtime comparisons fair, all benchmark runs were generated by the same experiment driver and executed on the same machine configuration, using the available CPU threads. The next chapter presents the measured runtime and route-quality results for each method.

\section{Benchmark summary}

Table~\ref{tab:benchmark-summary} collects the main numerical results for the compared methods. The shortest executable replay was obtained both by A* and by the tiny DQN trained with the Levenberg--Marquardt optimizer, while the Adam-trained variants provided a useful compromise between runtime and solution quality. In contrast, the EKF-based learned solvers and the purely route-based stochastic optimizers remained clearly above the optimal replay length on this benchmark.

\begin{table}[h]
    \centering
    \caption{Summary of benchmark results on the fixed wrapped $20 \times 20$ board. Replay length is measured on the exported executable trajectory; $\Delta$opt denotes the number of replay steps above the A* reference.}
    \label{tab:benchmark-summary}
    \begin{tabular}{lrrr}
        \toprule
        Method & Replay steps & $\Delta$opt & Wall time [s] \\
        \midrule
        A* & 71 & 0 & 0.10 \\
        Genetic & 95 & 24 & 5.09 \\
        Simulated annealing & 121 & 50 & 0.25 \\
        Metropolis & 97 & 26 & 0.26 \\
        Tiny DQN + EKF & 89 & 18 & 91.61 \\
        Tiny DQN + LM & 71 & 0 & 75.81 \\
        Tiny DQN + Adam & 77 & 6 & 7.44 \\
        Kohonen + EKF & 103 & 32 & 10.55 \\
        Kohonen + LM & 85 & 14 & 4.73 \\
        Kohonen + Adam & 86 & 15 & 1.73 \\
        \bottomrule
    \end{tabular}
\end{table}

\end{document}
